\section*{Solutions for Exercises from Section 15}
\begin{solution}\label{sol:2.15.1}\ref{ex:2.15.1}
    For $(1)$: 
    \begin{table}[H]
    \centering
    \begin{tabular}{ccccc}
    \toprule
    $P$ & $Q$ & $R$ & $Q \Rightarrow R$ & $P \lor (Q \Rightarrow R)$ \\
    \midrule
    T & T & T & T & T \\
    T & T & F & F & T \\
    T & F & T & T & T \\
    T & F & F & T & T \\
    F & T & T & T & T \\
    F & T & F & F & \textbf{F} \\
    F & F & T & T & T \\
    F & F & F & T & T \\
    \bottomrule
    \end{tabular}
    \end{table}

    $(2)$:
    \begin{table}[H]
        \centering
        \begin{tabular}{ccc}
        \toprule
        $Q$ & $R$ & $(Q\lor R)\Leftrightarrow (R\land Q)$\\
        \midrule
        T & T & T\\
        T & F & F\\
        F & T & F\\
        F & F & T\\
        \bottomrule
        \end{tabular}
    \end{table}    
   
    $(3)$:
    \begin{table}[H]
        \centering
        \begin{tabular}{ccc}
            \toprule
            $P$ & $Q$ & $\neg(P\Rightarrow Q)$\\
            \midrule
            T & T & F \\
            T & F & T \\
            F & T & F \\
            F & F & F\\
            \bottomrule
        \end{tabular}
    \end{table}

    $(4)$:
    \begin{table}[H]
        \centering
        \begin{tabular}{ccc}
            \toprule
            $P$ & $Q$ & $\neg(P\lor Q)\lor(\neg P)$\\
            \midrule
            T & T & F\\
            T & F & F\\
            F & T & T\\
            F & F & T\\
            \bottomrule
        \end{tabular}
    \end{table}

    $(5)$:
        \begin{table}[H]
        \centering
        \begin{tabular}{ccc}
            \toprule
            $P$ & $Q$ & $(P\land\neg P)\lor Q$\\
            \midrule
            T & T & T\\
            T & F & F\\
            F & T & T\\
            F & F & F\\
            \bottomrule
        \end{tabular}
    \end{table}

    $(6)$: 
        \begin{table}[H]
        \centering
        \begin{tabular}{ccc}
            \toprule
            $P$ & $Q$ & $(P\land \neg P)\land Q$\\
            \midrule
            T & T & F\\
            T & F & F\\
            F & T & F\\
            F & F & F\\
            \bottomrule
        \end{tabular}
    \end{table}

    $(7)$:
        \begin{table}[H]
        \centering
        \begin{tabular}{ccc}
            \toprule
            $P$ & $Q$ & $\neg(P\land \neg P)\Rightarrow Q$\\
            \midrule
            T & T & T\\
            T & F & T\\
            F & T & T\\
            F & F & T\\
            \bottomrule
        \end{tabular}
    \end{table}

    $(8)$:
    \begin{table}[H]
        \centering
        \begin{tabular}{cccc}
            \toprule
            $P$ & $Q$ & $R$ & $P\lor(Q\land\neg R)$\\
            \midrule
            T & T & T & T\\
            T & T & F & T\\
            T & F & T & T\\
            T & F & F & T\\
            F & T & T & F\\
            F & T & F & T\\
            F & F & T & F\\
            F & F & F & F\\
            \bottomrule
        \end{tabular}
    \end{table}

    $(9)$:
        \begin{table}[H]
        \centering
        \begin{tabular}{ccc}
            \toprule
            $P$ & $Q$ & $\neg(P\lor\neg Q)$\\
            \midrule
            T & T & T\\
            T & F & F\\
            F & T & F\\
            F & F & F\\
            \bottomrule
        \end{tabular}
    \end{table}    
\end{solution}

\begin{solution}\label{sol:2.15.2}\ref{ex:2.15.2}
    First we will find it without a truth table, and then we will confirm our answer by using a truth table.\newline
    First, we need to find the global operator, in our case 
    $$
    ((P\land Q)\lor R)\Rightarrow (R\lor S)
    $$
    this is the `$\Rightarrow$'. So, we know that if we have a true hypothesis and a false conclusion, then the conditional will be false. Therefore, we require that $P=T,Q=T, R=F, S=F$. Now to check with our truth table:
    \begin{table}[H]
        \centering
        \begin{tabular}{ccccc}
            \toprule
            $P$ & $Q$ & $R$ & $S$ & $((P\land Q)\lor R)\Rightarrow (R\lor S)$\\
            T & T & T & T & T\\
            T & T & T & F & T\\
            T & T & F & T & T\\
            T & T & F & F & \textbf{F}!\\
            T & F & T & T & T\\
            T & F & T & F & T\\
            T & F & F & T & T\\
            T & F & F & F & T\\
            F & T & T & T & T\\
            F & T & T & F & T\\
            F & T & F & T & T\\
            F & T & F & F & T\\
            F & F & T & T & T\\
            F & F & T & F & T\\
            F & F & F & T & T\\
            F & F & F & F & T\\
            \bottomrule
        \end{tabular}
    \end{table}
    So, everything checks out.
\end{solution}

\begin{solution}\label{sol:2.15.3}\ref{ex:2.15.3}
    We will follow the same approach. The formula is 
    $$
    (R\Rightarrow S)\Leftrightarrow (P\land Q)
    $$
    and we are supposing that it is true. So we learned that the biconditional requires both sides to have the same truth value. Then if $R=T$ and $P,Q=T$. It follows that the biconditional is true. Let us examine the truth table:
    \begin{table}[H]
        \centering
        \begin{tabular}{ccccc}
            \toprule
            $P$ & $Q$ & $R$ & $S$ & $((P\land Q)\lor R)\Rightarrow (R\lor S)$\\
            T & T & T & T & T\\
            T & T & T & F & F (because of $S$)\\
            T & T & F & T & T\\
            T & T & F & F & T\\
            T & F & T & T & F (because of $Q$)\\
            T & F & T & F & T\\
            T & F & F & T & F (because of $Q$)\\
            T & F & F & F & F (because of $Q$)\\
            F & T & T & T & F (because of $P$)\\
            F & T & T & F & T\\
            F & T & F & T & F (because of $P$)\\
            F & T & F & F & F (because of $P$)\\
            F & F & T & T & F (because of $P,Q$)\\
            F & F & T & F & T\\
            F & F & F & T & F\\
            F & F & F & F & F\\
            \bottomrule
        \end{tabular}
    \end{table}
\end{solution}